\section{Pendahuluan}

Efek Warburg --- fenomena di mana sel tumbuh-cepat memfermentasi glukosa
menjadi laktat atau etanol meskipun oksigen tersedia --- merupakan salah
satu teka-teki terlama dalam biologi sel \cite{warburg1956}.
Penjelasan yang dominan selama dua dekade terakhir berpusat pada
\textit{efisiensi proteom}: sel memiliki anggaran protein terbatas, dan
enzim glikolitik diduga menghasilkan ATP lebih cepat per gram protein
dibandingkan rantai pernapasan yang besar \cite{basan2015, pfeiffer2001}.

Pada tahun 2024, dua makalah tinjauan-sejawat menyimpulkan hal
berlawanan. Shen dkk.\ \cite{shen2024} menggunakan fluks metabolik
$^{13}$C-MFA dan proteomik kuantitatif pada ragi, sel T, dan 60 lini
sel kanker NCI60, melaporkan bahwa respirasi \textit{beberapa kali lipat
lebih efisien-proteom} daripada glikolisis. Sebaliknya, Kukurugya dkk.\
\cite{kukurugya2024} menyajikan model 5-parameter analitik dan data
eksperimental yang menunjukkan glikolisis menghasilkan ATP $0{,}54\times$
hingga $3{,}1\times$ lebih cepat per miligram enzim jalur dibandingkan
respirasi, bergantung pada organisme.

Wang \cite{wang2025} telah mendamaikan keduanya dengan memperkenalkan
model baru yang menggabungkan alokasi protein dan heterogenitas sel,
menunjukkan bahwa efisiensi respirasi dan fermentasi \textit{berpotongan}
menurut kualitas nutrisi. Namun, belum ada yang melakukan
\textbf{perbandingan kepala-ke-kepala} terhadap kedua model asli untuk
melokalisasi sumber perbedaan --- apakah perbedaan itu berasal dari
(i)~besaran yang diukur, (ii)~definisi ``protein jalur'',
(iii)~nilai parameter, atau (iv)~struktur model --- serta belum ada yang
menanyakan \textbf{pengukuran tunggal mana yang akan membalikkan kesimpulan}.

Paper ini memberikan tiga kontribusi:
\begin{enumerate}
    \item Implementasi ulang kedua model 2024 dalam satu basis kode
          dengan antarmuka bersama dan gerbang validasi terhadap gambar terbit.
    \item Tiga eksperimen diagnostik: overlay sumbu-bersama (E1),
          kapasitas versus realisasi (E2), dan analisis identifiabilitas
          Sobol/Morris (E3) yang menamai pengukuran penuntas.
    \item Kuantifikasi ketidakpastian (bootstrap 90\% CI) pada titik-balik
          verdict, dekomposisi atribusi enzim, dan perbandingan terhadap
          mekanisme penyilangan Wang~2025.
\end{enumerate}
